\chapter{Conclusiones y líneas futuras}

\section{Conclusiones}

El punto de partida del trabajo era una ciudad preparada para la conducción del usuario, pero sin un sistema de tráfico capaz de recorrerla de forma autónoma. Las primeras pruebas confirmaron que trasladar soluciones habituales de navegación de personajes a un coche con físicas producía recorridos rígidos y problemas en los cruces. La decisión de sustituir las rutas cerradas por percepción local y una política neuronal permitió avanzar, aunque el resultado no debe interpretarse como un controlador ya resuelto.

La arquitectura implementada conecta once sensores, velocidad, giro y contexto de señalización con una red de 20 entradas, 16 neuronas ocultas y dos salidas. Sus 352 pesos se validan antes de cada inferencia y las consignas se aplican de forma suavizada al componente físico. El gestor evolutivo conserva por separado el modelo histórico y el de sesión, distribuye mutaciones de distinta intensidad y normaliza la selección por punto de aparición. Esta integración funciona dentro del escenario real del simulador y no en una prueba aislada, que era una de las dificultades principales del proyecto.

El desarrollo también ha puesto de manifiesto que la función de fitness y el registro de datos forman parte de la solución, no son tareas posteriores. Una demostración visual podía parecer correcta y ocultar estancamientos, penalizaciones mal equilibradas o resultados concentrados en unos pocos puntos del mapa. La exportación por vehículo permitió revisar esas situaciones y motivó cambios concretos en el fitness, en los bordes de intersección y en la separación entre entrenamiento y test.

Los cuatro tests del 18 de junio reúnen 1134 evaluaciones y alcanzan un 57.50\% de acierto operativo, con un intervalo de confianza del 95\% entre el 54.60\% y el 60.34\%. Las colisiones descienden al 40.04\% y los fallos legales representan el 2.29\%. La mejora frente al bloque del día 8 es notable, pero el rendimiento sigue dependiendo del punto de aparición y las rotondas acumulan un saldo de fitness negativo. El entrenamiento asociado tampoco presenta convergencia clara: aumenta el tiempo de vida mientras desciende el fitness medio entre el inicio y el final de la sesión.

Las ejecuciones posteriores refuerzan esa cautela. Después de modificar la escala del fitness, la lógica de parada y el tratamiento de intersecciones, la serie alterna recuperaciones parciales con regresiones marcadas por \texttt{STOP\_VIOLATION}. Algunas sesiones aisladas obtienen resultados aceptables, pero los bloques amplios vuelven a mostrar inestabilidad y diferencias importantes entre causas de fallo. Las últimas pruebas aportan métricas más finas sobre \texttt{STOP\_VIOLATION\_TIMEOUT}, cedas liberados, bloqueo en cruces, solapes entre coches y espera en colas dinámicas de ceda. La segunda sesión \texttt{FreeRun} del 3 de julio elimina \texttt{LAZY} como causa terminal, pero desplaza el fallo dominante hacia \texttt{NAV\_VIOLATION} y \texttt{STOP\_VIOLATION\_EXIT}. Esta evolución se interpreta como una regresión detectada y mejor instrumentada, no como una nueva referencia. El resultado muestra el valor práctico de conservar configuraciones y de contrastar cada cambio con más de una métrica.

Por tanto, se ha obtenido una base funcional para generar, entrenar y medir tráfico autónomo, pero todavía no un sistema suficientemente robusto para una sesión clínica. La utilidad inmediata del trabajo reside en que los fallos ya pueden localizarse y compararse con un protocolo común. Las siguientes iteraciones deberán demostrar que las mejoras se mantienen en intersecciones y puntos de inicio distintos antes de aumentar la densidad de tráfico o incorporar nuevos tipos de agentes.

\clearpage
\section{Aportaciones realizadas}

La primera aportación es el controlador neuronal integrado en el vehículo de Unreal Engine. El trabajo incluye la construcción del vector de 20 entradas, la inferencia de la red y la aplicación suavizada de las dos salidas al sistema Chaos Vehicle. Los sensores no se limitan a detectar obstáculos: en las intersecciones filtran los bordes según la pareja de trigger y maniobra elegida y combinan esa geometría con el estado de semáforos, stops y cedas.

La segunda aportación es el proceso de aprendizaje dentro del propio simulador. El gestor crea poblaciones, conserva sin mutar las dos semillas principales, aplica mutaciones de intensidad variable y guarda por separado el mejor modelo histórico y el mejor de sesión. La versión final calcula la tasa de mutación a partir del historial de generaciones sin mejora y del número total de pesos, en lugar de depender únicamente de porcentajes fijos. La referencia de cada spawn decae con el historial para evitar máximos permanentes, y el uso de puntos y opciones de trayectoria distintos en test permite evaluar el modelo fuera del conjunto que condicionó su entrenamiento.

La tercera aportación es la infraestructura de evaluación. Los Blueprints exportan el resultado de cada vehículo y los resúmenes de generación; posteriormente, los scripts comprueban la calidad de los CSV, agrupan razones de muerte y producen comparaciones por sesión, spawn, familia de mutación, solapes entre vehículos, bloqueo en cruce, validación y cola de ceda, y salida de rotonda. Esta conexión entre simulación y análisis permite justificar una modificación con datos y reconocer una regresión aunque el cambio local pareciera razonable.

Otra aportación relevante es la transición hacia tráfico físicamente concurrente. La incorporación de coches no fantasma, detección de solapes, penalización real por distancia de seguridad, zonas de respawn libres y simulación continua con \texttt{FreeRun} acerca el sistema al uso previsto: generar tráfico dinámico alrededor del usuario, no únicamente evaluar coches aislados contra el mapa. Esta mejora aumenta el realismo y también revela nuevas fuentes de fallo, como colas, bloqueos y conflictos entre vehículos.

Finalmente, las actas, las issues de GitLab y las descripciones fechadas conservan la evolución técnica del proyecto. Gracias a ellas puede reconstruirse por qué se abandonó \textit{Path Finding}, cuándo se reformuló el fitness o cómo se introdujeron los bordes dependientes de la trayectoria. El registro depurado de 40 issues cerradas aporta una capa adicional de trazabilidad entre reuniones, implementación y resultados. Esto resulta especialmente útil en un desarrollo compartido que deberá continuar después del TFG\@.

\section{Problemas encontrados}

El primer problema fue adaptar la navegación de Unreal Engine a un vehículo con inercia y radio de giro. NavMesh y \textit{Path Finding} encontraban un destino, pero no producían una conducción natural. Abandonar esa vía obligó a trabajar conjuntamente con Chaos Vehicle, el control mediante Blueprints y los sensores de calzada.

Las intersecciones concentraron buena parte de los fallos posteriores. Un mismo cruce debía ofrecer límites distintos según la salida elegida, y esos límites tenían que coincidir con el trigger que había originado la decisión. La asociación explícita de parejas formadas por dirección y trigger resolvió parte del problema y permitió reutilizar bordes comunes. Aun así, \texttt{BP\_IntersectionBorder221} concentra 94 colisiones en el bloque de referencia, y las rotondas fallan con mayor frecuencia en recorridos que alcanzan salidas posteriores.

El ajuste del fitness fue otra dificultad importante. Premiar avance y supervivencia generó estrategias no previstas, entre ellas permanecer detenido o compensar una colisión con recompensas acumuladas. En rotondas ocurrió el problema contrario: la penalización de aceleración superó ampliamente las bonificaciones. La corrección posterior de la escala coincidió con una regresión masiva de stops, y las pruebas de seguimiento confirmaron que una ejecución favorable aislada no basta para cerrar el problema. La instrumentación aumentó la complejidad de los datos, pero evitó continuar ajustando el modelo únicamente a partir de la impresión visual.

Por último, cada entrenamiento exige ejecutar física y renderizado para una población completa. Este coste limita la velocidad con la que pueden probarse cambios y hace especialmente caro descubrir tarde un error de configuración. La persistencia de modelos, la validación de los CSV y los bloques de test repetidos reducen ese riesgo, aunque sigue pendiente un modo de entrenamiento más rápido.

\section{Líneas futuras}

La siguiente iteración debería comenzar por los puntos 3, 36 y 42 y por tres actores concretos del bloque de referencia: el borde de intersección 221 y los volúmenes de navegación 43 y 35. Además, la sesión \texttt{FreeRun} del 3 de julio señala nuevos focos de revisión en \texttt{BP\_SpawnPoint40} y en los triggers \texttt{IT58}, \texttt{IT55} e \texttt{IT54}. En cada caso conviene separar un error de geometría de un fallo del controlador: revisar orientación y bordes, repetir la primera maniobra con el modelo del día 18 y comparar distancia recorrida, tiempo detenido y causa de finalización. Los puntos 16, 9 y 25 deben mantenerse como control para detectar regresiones en recorridos que ahora funcionan mejor.

Antes de un nuevo entrenamiento largo debe validarse la corrección introducida en la lógica de stop el 6 de julio. La prueba adecuada consiste en mantener el modelo, el fitness restante y los spawns, variar solo esa lógica de validación y ejecutar varios bloques cortos de regresión. Después seguirá siendo necesario distinguir una supervivencia útil de un coche inmóvil ante una señal mediante progreso de ruta, duración continua del atasco, tiempo bloqueado en cruce y subtipo de infracción. Los tests deberán conservar tamaño, modelo y conjunto de spawns para que los intervalos de confianza sean comparables.

También debe equilibrarse el término de rotonda. Las 31 completaciones confirmadas frente a 443 entradas y el saldo neto negativo aconsejan validar por separado la penalización de entrada, la de aceleración y el bonus de salida. Las salidas segunda y tercera requieren pruebas dirigidas, porque acumulan una proporción de muertes superior a la primera.

Para reducir el tiempo entre iteraciones se plantea un modo sin renderizado o con la simulación desacoplada de la presentación. Esta mejora permitiría entrenar más generaciones y reservar la ejecución completa para validar física, señalización y comportamiento visual. También sería conveniente automatizar pruebas sobre intersecciones concretas antes de lanzar una sesión larga.

\Needspace{9\baselineskip}
Otra línea inmediata es integrar mejor la configuración del experimento con \textit{VictoryPlugin}. El objetivo sería exportar e importar perfiles de ejecución con variables clave como tamaño de población, umbrales de test, tiempos de vida, modo train/test, activación de \texttt{FreeRun}, selección de spawns y parámetros de mutación. Con ello podrían ofrecerse perfiles comprensibles para distintos usos: un \textit{Modo Ligero}, con menor población y ejecución más rápida para depuración, y un \textit{Modo Completo}, con más spawns y métricas extensas para validación. Esta configuración reduciría errores manuales antes de entrenamientos largos.

También debe evolucionar el sistema de puntos de aparición. Actualmente los \texttt{BP\_SpawnPoint} permiten separar train/test y comprobar si la zona está libre, pero siguen dependiendo de colocación manual. Una mejora consistiría en generar spawns autónomos, espaciados y orientados de forma coherente con el carril, ajustando automáticamente la cantidad a la población elegida. Esto permitiría que el mapa se adapte al experimento, evitaría concentraciones de coches en zonas concretas y facilitaría comparar sesiones con distintas densidades de tráfico.

En accesibilidad, el simulador debería incorporar ayudas para usuarios invidentes o con baja visión. Además de mejorar contraste, tamaño de elementos y claridad de instrucciones, se pueden estudiar estímulos hápticos mediante vibraciones y señales acústicas avanzadas, por ejemplo pitidos diferenciados para proximidad, semáforo, salida de carril o riesgo de colisión. La compatibilidad con volantes y pedales adaptados también es prioritaria si el simulador se orienta a rehabilitación y valoración funcional. Estas mejoras deberán diseñarse con profesionales y usuarios, no solo desde criterios técnicos.

Por último, el comportamiento de conducción avanzada queda abierto para fases posteriores. Una vez estabilizada la circulación básica, el tráfico autónomo debería incorporar cambio de carril, adelantamientos, incorporación a vías, gestión de prioridad entre varios coches y conducción anticipativa ante retenciones. Estas lógicas requieren que la red o el gestor dispongan de contexto adicional, porque no basta con reaccionar al obstáculo frontal: el vehículo debe decidir cuándo esperar, cuándo ceder, cuándo rebasar y cómo coordinarse con otros agentes.

Los perfiles de conductor normal, inseguro y agresivo, la interacción entre vehículos y la incorporación de peatones o bicicletas deben abordarse cuando la base vehicular sea estable. Añadir agentes antes de resolver las colisiones actuales dificultaría atribuir cada fallo a su causa. Finalmente, cualquier uso con pacientes requerirá definir con ASPAYM un protocolo de sesión, criterios de accesibilidad e indicadores comprensibles para los profesionales.

\chapter*{Agradecimientos}
\phantomsection%
\addcontentsline{toc}{chapter}{Agradecimientos}
